\chapter{Auxiliary circuits}
\label{ch:aux}

\section{Gate drive}

\subsection{Requirements}

The peak gate current follows from the gate charge and the total gate-loop
resistance:
%
\begin{equation}
  I_{g,\mathrm{pk}} = \frac{V_\mathrm{DD} - V_\mathrm{EE}}
                           {R_{g,\mathrm{on}} + R_{G(\mathrm{int})} + R_\mathrm{drv}}
  = \frac{\SI{20}{\volt}}{\SI{5}{\ohm} + \SI{2.6}{\ohm} + \SI{2.5}{\ohm}}
  = \SI{1.98}{\ampere},
\end{equation}
%
and the average gate-drive power per channel from the gate charge:
%
\begin{equation}
  P_g = Q_g \, (V_\mathrm{DD}-V_\mathrm{EE}) \, \fsw
      = \SI{223}{\nano\coulomb} \cdot \SI{20}{\volt} \cdot \SI{20}{\kilo\hertz}
      = \SI{89.2}{\milli\watt}.
\end{equation}
%
The \SI{20}{\volt} swing is $+15$ to $-5$~V, the choice justified in
\cref{sec:negbias}.

\subsection{Driver selection}

\begin{table}[htbp]
\centering
\caption{Isolated gate driver alternatives.}
\label{tab:driver}
\small
\begin{tabularx}{\textwidth}{lYYY}
\toprule
 & \textbf{\driver{}} & Basic isolated driver & Optocoupler driver \\
\midrule
Isolation & \textbf{reinforced, \SI{2121}{\volt} DC} & basic & reinforced \\
CMTI & \textbf{\SI{150}{\volt\per\nano\second}} & \SI{50}{}--\SI{100}{\volt\per\nano\second} & \SI{25}{}--\SI{50}{\volt\per\nano\second} \\
Desaturation detection & \textbf{integrated} & external & some \\
Active Miller clamp & \textbf{integrated} & external & rare \\
Soft turn-off & \textbf{integrated} & no & no \\
Propagation delay & \SI{90}{\nano\second} & similar & \SI{150}{}--\SI{300}{\nano\second} \\
Ageing & \textbf{capacitive, stable} & --- & LED degrades \\
\bottomrule
\end{tabularx}
\end{table}

The \driver{} is selected~\cite{ti_ucc21750}; \cref{fig:gatechannel} shows
one channel as built. The decisive parameter is the
common-mode transient immunity. In the double-pulse simulation at the
overvoltage and overload thresholds (\cref{sec:kelvin}) the switching node
moves at \SI{31}{\volt\per\nano\second} between \SI{10}{} and
\SI{90}{\percent} of the steeper edge, and at \SI{70}{\volt\per\nano\second}
over the steepest nanosecond, ringing included. A driver rated
\SI{50}{\volt\per\nano\second} would be exceeded; the
\SI{150}{\volt\per\nano\second} rating leaves a factor of 2.1 on the
steepest value.

Optocoupler-based drivers are rejected on two grounds: their propagation
delay is two to three times larger, which directly enlarges the dead-time
budget of \cref{tab:deadtime}, and the LED degrades with time and
temperature, which is a poor property in an automotive application.

\begin{figure}[htbp]
\centering
% Un canale di gate: driver isolato, resistenze separate di accensione e spegnimento, sorgente
% Kelvin, rete di desaturazione con zener, alimentazione isolata +15/-5 V.
\begin{circuitikz}[font=\footnotesize, >=Stealth]
  % driver isolato (pin disposti per non incrociare i collegamenti)
  \draw[thick, fill=fblu!6] (0,0) rectangle (3.0,5.3);
  \node[font=\footnotesize\bfseries] at (1.5,5.6) {isolated driver};
  \draw[very thick, dashed, fgrigio] (1.0,-0.25) -- (1.0,5.5);
  \node[font=\scriptsize, rotate=90, fgrigio] at (0.7,2.6) {reinforced barrier};
  \foreach \y/\t in {4.2/PWM, 3.4/enable, 2.6/fault, 1.8/ready}
    \draw[<->, fblu] (-1.0,\y) node[left]{\t} -- (0,\y);
  \foreach \y/\t in {4.9/DESAT, 3.5/OUTH, 2.7/OUTL, 1.7/COM, 1.0/VDD, 0.4/VEE}
    \node[anchor=east, font=\scriptsize] at (2.95,\y) {\t};
  % resistenze di gate separate
  \draw (3.0,3.5) to[R, l={$R_{g,\mathrm{on}}$, \SI{5}{\ohm}}] (5.6,3.5) -- (6.3,3.5) -- (6.3,3.1) coordinate (g);
  \draw (3.0,2.7) to[R, l_={$R_{g,\mathrm{off}}$, \SI{12.5}{\ohm}}] (5.6,2.7) -- (6.3,2.7) -- (g);
  \draw (g) -- (7.4,3.1) node[nmos, bodydiode, anchor=G] (Q) {};
  \draw (Q.D) -- ++(0,1.8) node[above]{drain};
  \draw (Q.S) -- ++(0,-0.9) node[below]{power source};
  % sorgente Kelvin: il ritorno del driver va al source sul die, non sul percorso di potenza
  \draw (3.0,1.7) -- ($(Q.S |- 0,1.7)+(-0.35,0)$) |- ($(Q.S)+(0,-0.3)$);
  \node[font=\scriptsize, anchor=north] at ($(Q.S |- 0,1.7)+(-1.4,0)$) {Kelvin source};
  % desaturazione: resistenza, condensatore di blanking verso COM, zener, diodo HV al drain
  \draw (3.0,4.9) to[R, l={\SI{1}{\kilo\ohm}}, -*] (4.2,4.9) coordinate (dn);
  \draw (dn) to[C, l_={$C_\mathrm{blank}$}] (4.2,6.0) node[above, font=\scriptsize]{to COM};
  \draw (dn) -- (5.0,4.9) to[zDo, l_={\SI{4.3}{\volt} zener}] (6.4,4.9)
        to[D, l_={SiC diode}, -*] ($(Q.D |- 0,4.9)$);
  % alimentazione isolata del lato secondario
  \draw[thick, fill=fverde!8] (0.2,-1.7) rectangle (3.0,-0.7);
  \node[align=center, font=\scriptsize] at (1.6,-1.2) {isolated DC/DC\\$+\SI{15}{}/-\SI{5}{\volt}$};
  \draw[->] (3.0,-0.95) -| (3.5,1.0) -- (3.0,1.0);
  \draw[->] (3.0,-1.45) -| (3.8,0.4) -- (3.0,0.4);
\end{circuitikz}

\caption{One gate-drive channel. Separate turn-on and turn-off resistors, the
driver's return to the Kelvin source rather than to the power source, the
desaturation network with the \SI{4.3}{\volt} zener of
\cref{sec:protection}, and the isolated $+\SI{15}{}/-\SI{5}{\volt}$
supply.}
\label{fig:gatechannel}
\end{figure}

\subsection{Gate resistance and negative bias}

$R_{g,\mathrm{on}} = \SI{5}{\ohm}$ and $R_{g,\mathrm{off}} = \SI{12.5}{\ohm}$,
on the two separate outputs of the driver, with footprints allowing up to
\SI{25}{\ohm}. They are the smallest values that keep the die below \SI{1080}{\volt} at \SI{950}{\volt} and \SI{60}{\ampere} with
the \SI{20}{\nano\henry} layout budget (\cref{sec:kelvin}): $R_{g,\mathrm{off}}$
is set by the turn-off of the device itself, $R_{g,\mathrm{on}}$ by the
recovery of the opposite diode, and the separate outputs are what allow the
two to differ. The first version of this design
used \SI{2.5}{\ohm} on the argument that overshoot margin is cheaper to buy
from the layout; that argument holds for margin \emph{beyond} the
criterion, not for the criterion itself.

\subsection{Negative bias: why \texorpdfstring{$-\SI{5}{\volt}$}{-5 V}}
\label{sec:negbias}

A negative off-state bias is mandatory. With the gate
held at \SI{0}{\volt}, the Miller current injected through $C_\mathrm{gd}$
during the opposite device's turn-on raises the gate voltage by
%
\begin{equation}
  \Delta V_\mathrm{GS} = C_\mathrm{gd,eff}\,\frac{\mathrm{d}v}{\mathrm{d}t}\,R
  \left[1 - \exp\!\left(-\frac{t}{R\,C_\mathrm{iss}}\right)\right].
  \label{eq:miller}
\end{equation}

The datasheet quotes $-\SI{4}{\volt}$ as the \emph{typical operating}
value, not as a limit: the absolute maximum is
$-\SI{8}{\volt}$~\cite{wolfspeed_e3m0016}. This design uses
$-\SI{5}{\volt}$, which is what the standard isolated module delivers
without an extra regulator, and which leaves \SI{3}{\volt} to the absolute
maximum.

Since \cref{eq:miller} does not depend on the bias, the extra volt
translates one-for-one into margin. The equation is evaluated with
$C_\mathrm{gd,eff} = Q_\mathrm{gd}/V_\mathrm{DS} = \SI{81}{\pico\farad}$,
$C_\mathrm{iss} = \SI{6922}{\pico\farad}$ and the turn-on slope that the
double-pulse simulation gives for the opposite device,
\SI{31.4}{\volt\per\nano\second}; the closed-form slope from the gate current
would be only \SI{14.3}{\volt\per\nano\second}, and the steeper value is the
one used. Without the clamp the off device is held through
$R_{g,\mathrm{off}} + R_{G(\mathrm{int})}$ plus the driver's sink
resistance, \SI{15.4}{\ohm}. The same case was then simulated directly, holding the off device
at \SI{-5}{\volt} through its gate resistor, and reading $V_\mathrm{GS}$ at
the die:

\begin{table}[htbp]
\centering
\caption{Margin against spurious turn-on at \SI{950}{\volt},
\SI{60}{\ampere}. $V_\mathrm{GS(th)}$ is the \SI{1.8}{\volt} minimum at
\SI{25}{\celsius}~\cite{wolfspeed_e3m0016}.}
\label{tab:miller}
\begin{tabular}{lrrr}
\toprule
Configuration & \Cref{eq:miller} & Margin & Simulated, at the die \\
\midrule
$-\SI{4}{\volt}$, no clamp   & $+\SI{4.35}{\volt}$ & $-\SI{2.55}{\volt}$ & --- \\
$-\SI{4}{\volt}$, with clamp & $+\SI{1.58}{\volt}$ & \SI{0.22}{\volt} & --- \\
$-\SI{5}{\volt}$, no clamp   & $+\SI{3.35}{\volt}$ & $-\SI{1.55}{\volt}$ & $-\SI{1.33}{\volt}$ \\
\textbf{$-\SI{5}{\volt}$, with clamp} & $\mathbf{+\SI{0.58}{\volt}}$ & \textbf{\SI{1.22}{\volt}} & $\mathbf{-\SI{2.98}{\volt}}$ \\
\bottomrule
\end{tabular}
\end{table}

The clamp is still required. The closed form is a bound, not a
prediction: it charges the whole Miller current into the external resistance
for the whole edge, and the simulation, which resolves the internal gate
network, keeps the die \SIrange{4}{5}{\volt} lower. But the bound is the
number that can be defended without trusting a model's gate network, and by
that number $-\SI{5}{\volt}$ alone fails and $-\SI{5}{\volt}$ with the clamp
passes with \SI{1.22}{\volt}. The integrated active Miller clamp removes the
dependence on $R_{g,\mathrm{off}}$ by shorting the gate through a separate
low-impedance path; with \SI{12.5}{\ohm} of $R_{g,\mathrm{off}}$ it is what
makes the margin exist.

One effect runs the other way: the body
diode forward voltage is characterised at
$V_\mathrm{GS} = -\SI{4}{\volt}$~\cite{wolfspeed_e3m0016}, so a more negative
gate raises it slightly and with it the dead-time loss of
\cref{sec:deadtime-loss}. The effect is a few hundredths of a watt, under
\SI{0.05}{\percent} of the total, but it is the reason not to go to
$-\SI{8}{\volt}$ merely because the margin allows it.

\section{Protection}
\label{sec:protection}

\subsection{Overcurrent: two mechanisms, two roles}

\begin{table}[htbp]
\centering
\caption{Protection assignment.}
\label{tab:protection}
\small
\begin{tabularx}{\textwidth}{lllY}
\toprule
Mechanism & Threshold & Response time & Covers \\
\midrule
Desaturation, with \SI{4.3}{\volt} zener & \SIrange{143}{259}{\ampere} & \SI{0.69}{}/\SI{0.83}{\micro\second}$^\ast$ & hard short, shoot-through \\
Shunt + comparator & $\approx\SI{60}{\ampere}$ & \SI{3}{}--\SI{4}{\micro\second} & overload, ground fault, stall \\
\bottomrule
\multicolumn{4}{l}{\footnotesize $^\ast$ typical/worst case, counted as $t_\mathrm{SC}$ is measured (\cref{tab:scchain}).}
\end{tabularx}
\end{table}

\begin{figure}[htbp]
\centering
% Scale di tempo: ogni protezione e' assegnata al guasto che riesce a raggiungere in tempo.
% Asse logaritmico; ogni barra va dall'inizio dell'evento al suo tempo (capitoli 7 e 8).
\begin{tikzpicture}
  \begin{axis}[
    width=0.74\textwidth, height=6.2cm, xmode=log, clip=false,
    xmin=1e-7, xmax=1e3, ymin=0.4, ymax=7.6,
    ytick={1,...,7},
    yticklabels={{soft turn-off of a fault}, {desaturation: hard short}, {short-circuit withstand},
                  {shunt and comparator: overload}, {active discharge to \SI{60}{\volt}},
                  {thermistor: cold-plate time constant}, {passive bleeder to \SI{60}{\volt}}},
    xlabel={time [s]}, xmajorgrids, grid style={black!10},
    tick label style={font=\scriptsize}, label style={font=\footnotesize}]
    % barre: dall'inizio dell'evento al tempo caratteristico
    \foreach \b/\r/\c in {2.5e-7/1/farancio, 0.83e-6/2/farancio, 1.66e-6/3/fgrigio, 3.7e-6/4/fblu,
                          0.54/5/fblu, 10/6/fverde, 121/7/fgrigio} {
      \edef\tmp{\noexpand\fill[\c!70] (axis cs:1e-7,\r-0.28) rectangle (axis cs:\b,\r+0.28);}\tmp
    }
    \node[font=\scriptsize, anchor=west] at (axis cs:2.8e-7,1) {about \SI{250}{\nano\second}};
    \node[font=\scriptsize, anchor=west] at (axis cs:0.95e-6,2) {\SI{0.69}{}--\SI{0.83}{\micro\second}};
    \node[font=\scriptsize, anchor=west] at (axis cs:1.9e-6,3) {\SI{0.83}{\micro\second} needed, \SI{1.66}{} with margin};
    \node[font=\scriptsize, anchor=west] at (axis cs:4.2e-6,4) {\SI{3.2}{}--\SI{3.7}{\micro\second}};
    \node[font=\scriptsize, anchor=west] at (axis cs:0.62,5) {\SI{0.54}{\second}};
    \node[font=\scriptsize, anchor=west] at (axis cs:12,6) {$\tau \approx \SI{10}{\second}$};
    \node[font=\scriptsize, anchor=east, text=white] at (axis cs:100,7) {\SI{121}{\second}, only without \SI{12}{\volt}};
    % requisito di scarica
    \draw[dashed, farancio, thick] (axis cs:5,4.5) -- (axis cs:5,7.6);
    \node[font=\scriptsize, farancio, anchor=south] at (axis cs:5,7.6) {\SI{5}{\second} required};
  \end{axis}
\end{tikzpicture}

\caption{Time scales of the protections and of the events they answer. Each layer
covers the faults it can reach in time: desaturation the hard short, inside
the withstand time; the shunt chain the overload; the thermistor the loss
of cooling; the active discharge the \SI{5}{\second} requirement, which the
passive bleeder alone misses.}
\label{fig:protecttime}
\end{figure}

\Cref{fig:protecttime} places every protection on the time scale of the
fault it answers. Desaturation detection on a SiC MOSFET is harder than on
an IGBT, and the
report should say so. The method compares the on-state drain voltage with a
threshold; an IGBT in short circuit presents several volts of
$V_\mathrm{CE(sat)}$, whereas a SiC MOSFET presents only $I\cdot\Rdson$. The
driver threshold sits at \SI{9.15}{\volt}~\cite{ti_ucc21750}; less the
blocking diode drop, the drain trips at about \SI{8.45}{\volt}.

\begin{finding}{With a \SI{16}{\milli\ohm} device, the native threshold trips beyond the pulsed rating}
$8.45 / \Rdson$ is \SI{291}{\ampere} with the die at \SI{175}{\celsius}
and \SI{528}{\ampere} at \SI{25}{\celsius}, against a pulsed drain rating of
\SI{321}{\ampere}~\cite{wolfspeed_e3m0016}. A cold device with a slowly
rising fault current (a short through a motor winding rather than across the
bridge) would reach its rating before the protection saw it. The threshold
that was comfortable with a \SI{40}{\milli\ohm} device is wrong with a
\SI{16}{\milli\ohm} one.
\end{finding}

A \SI{4.3}{\volt} zener in series with the blocking diode lowers the drain
threshold to about \SI{4.15}{\volt}: \SI{143}{\ampere} hot, \SI{259}{\ampere}
at \SI{25}{\celsius} with a typical device and \SI{189}{\ampere} with a
maximum-resistance one, all below the pulsed rating. Normal operation stays
well clear of it: at the \SI{42.4}{\ampere} rated peak a hot device shows
\SI{1.23}{\volt}, a factor of 3.4, and at the \SI{60}{\ampere} overload trip
\SI{1.74}{\volt}, a factor of 2.4. Those are typical figures. Stacking the
tolerances (the lowest detection threshold, \SI{8.5}{\volt}; a blocking
diode at \SI{0.8}{\volt}; the zener at $+\SI{5}{\percent}$; a hot device at
maximum resistance, taken as the typical hot value scaled by the
\num{22}/\num{16} ratio the datasheet gives at \SI{25}{\celsius}), the trip
point falls to \SI{3.19}{\volt} and the \SI{60}{\ampere} operating point
rises to \SI{2.39}{\volt}: a factor of 1.33. That factor, not 2.4, is what
limits any further lowering of the threshold (\cref{sec:planb}). One zener
per channel is the whole cost.

The shunt chain is deliberately \emph{not} assigned to short-circuit
protection, because it cannot reach it. The \isoamp{} alone has a
propagation delay of \SI{2.1}{\micro\second} maximum at
\SI{50}{\percent}--\SI{50}{\percent} and \SI{3.0}{\micro\second} to
\SI{90}{\percent}~\cite{ti_amc1302}, set by its sigma-delta modulator; with
the comparator and the soft turn-off the chain totals
\SI{3.2}{}--\SI{3.7}{\micro\second}. That is fine for an overload, which
develops over milliseconds, and useless against a short.

\subsection{Blanking}

The driver already provides \SI{200}{\nano\second} of leading-edge blanking
internally~\cite{ti_ucc21750}. The external capacitor adds only the
difference:
%
\begin{equation}
  t_\mathrm{add} = \frac{C_\mathrm{blank}\,V_\mathrm{DESAT}}{I_\mathrm{CHG}}
  = \frac{C_\mathrm{blank}\cdot\SI{9.15}{\volt}}{\SI{500}{\micro\ampere}} .
\end{equation}
%
With $C_\mathrm{blank} = \SI{10}{\pico\farad}$ the total is about
\SI{383}{\nano\second}, which covers
$t_\mathrm{d(on)} + t_r \approx \SI{81}{\nano\second}$ at
$R_{g,\mathrm{on}} = \SI{5}{\ohm}$~\cite{wolfspeed_e3m0016} with a wide
margin. Going
below \SI{10}{\pico\farad} is pointless: stray capacitance becomes
comparable with the component.

\subsection{Soft shutdown}

At \SI{60}{\ampere} and $R_{g,\mathrm{off}} = \SI{12.5}{\ohm}$ the die already
reaches \SI{1073.5}{\volt} (\cref{sec:kelvin}). Turning off a fault of
\SIrange{150}{260}{\ampere} through the same resistor would produce a
$\mathrm{d}i/\mathrm{d}t$ that, across the same loop, exceeds
\SI{1200}{\volt}: the protection would destroy the device it is protecting.
The driver's integrated soft turn-off pulls the gate down through a
\SI{400}{\milli\ampere} current source instead, lengthening the turn-off to
roughly \SI{250}{\nano\second} and keeping the overshoot bounded.

\subsection{Short-circuit withstand}
\label{sec:tsc}

The \device{} datasheet does not declare a short-circuit withstand
time. It is the one selection criterion the device does not demonstrate
(\cref{sec:deviceselection}), and the way to carry an unknown of this kind is
to turn it into a threshold: the value below which the protection no longer
protects.

\paragraph{The threshold.} The manufacturer measures $t_\mathrm{SC}$ from
\SI{50}{\percent} of the gate voltage at turn-on to \SI{50}{\percent} at
turn-off~\cite{wolfspeed_scan}, so the protection chain is counted the same
way, for the worst fault: a short already present at turn-on, which has to
wait out the blanking (\cref{tab:scchain}).

\begin{table}[htbp]
\centering
\caption{Desaturation chain, counted as $t_\mathrm{SC}$ is measured,
datasheet tolerances~\cite{ti_ucc21750}.}
\label{tab:scchain}
\begin{tabular}{lrr}
\toprule
Contribution & Typical & Worst case \\
\midrule
Internal leading-edge blanking & \SI{200}{\nano\second} & \SI{200}{\nano\second} \\
External blanking, $C_\mathrm{blank}V_\mathrm{DESAT}/I_\mathrm{CHG}$, \SI{10}{\pico\farad}
  & \SI{183}{\nano\second} & \SI{228}{\nano\second} \\
Detection to output low (deglitch included) & \SI{200}{\nano\second} & \SI{300}{\nano\second} \\
Soft turn-off, \SI{90}{\percent} to \SI{50}{\percent} of the gate swing & \SI{138}{\nano\second} & \SI{138}{\nano\second} \\
Less turn-on to \SI{50}{\percent} of the gate swing & $-\SI{36}{\nano\second}$ & $-\SI{36}{\nano\second}$ \\
\midrule
\textbf{Total} & \textbf{\SI{0.69}{\micro\second}} & \textbf{\SI{0.83}{\micro\second}} \\
\bottomrule
\end{tabular}
\end{table}

The design survives a short if $t_\mathrm{SC} \geq
\SI{0.83}{\micro\second}$, and has a factor of two at
\SI{1.66}{\micro\second}. The earlier figure of
\SIrange{0.65}{0.75}{\micro\second} was the typical chain; C5, fixed at
\SI{1.5}{\micro\second} before the tolerances were stacked, gives a factor
of 1.8 on the worst case.

\paragraph{What $t_\mathrm{SC}$ depends on.} During a short the device
sits in saturation with the whole bus across it, and it fails when the
energy $E_\mathrm{SC} \approx V_\mathrm{bus} I_\mathrm{sat} t$ reaches a
critical value set by the die: its area and its thermal mass, not the
circuit~\cite{wolfspeed_scan}. The circuit sets the other two factors.
$V_\mathrm{bus}$ is fixed. $I_\mathrm{sat}$ is set by the gate voltage,
roughly as $(V_\mathrm{GS} - V_\mathrm{th})^n$ with $n$ between 1 and 2.
This is where \SI{+15}{\volt} helps: it is the bottom of the range used
for SiC, where \SI{+18}{} and \SI{+20}{\volt} are common, and with
$V_\mathrm{th} = \SI{2.5}{\volt}$ the saturation current at \SI{+18}{\volt}
is \SIrange{24}{54}{\percent} higher, at \SI{+20}{\volt}
\SIrange{40}{96}{\percent} higher. For the same critical energy the
withstand time is shorter in the same proportion. A datasheet value
quoted at \SI{+18}{\volt} would be pessimistic for this design.

\paragraph{A class value.} The manufacturer has published a measurement on
the same die generation and resistance class: a \SI{16}{\milli\ohm} Gen~3
die in a half-bridge module survives \SI{2.9}{\micro\second} at
\SI{800}{\volt}, \SI{175}{\celsius} and \SI{+15}{\volt}, with a critical
energy of about \SI{1.4}{\joule}~\cite{wolfspeed_scan}. Scaled at constant
energy it becomes \SI{2.6}{\micro\second} at \SI{900}{\volt} and
\SI{2.4}{\micro\second} at \SI{950}{\volt}, 2.9 times the worst-case chain.
This is a class value, not a property of the \device{}: another
package, another die lot, an application note rather than a datasheet. It
says that the threshold is plausibly met with margin, not that it is met.
Peer-reviewed measurements on \SI{1.2}{\kilo\volt} devices show the same
strong dependence on bus voltage, for instance from \SI{28}{} to
\SI{7}{\micro\second} between \SI{300}{} and \SI{600}{\volt} on trench
devices at room temperature~\cite{zou2022sc}.

\paragraph{Plan B, sized now.}
\label{sec:planb}
If the bench returned a value between \SI{0.83}{} and
\SI{1.66}{\micro\second}, two changes are available without touching the
board, and one is not:
%
\begin{itemize}
  \item Blanking capacitor from \SI{10}{} to \SI{4.7}{\pico\farad}:
        the worst-case chain falls from \SI{0.83}{} to
        \SI{0.71}{\micro\second}. The internal \SI{200}{\nano\second}
        still covers the \SI{81}{\nano\second} turn-on by a factor of 2.5;
        below \SI{4.7}{\pico\farad} the component is comparable with the
        stray capacitance and filters nothing.
  \item Zener from \SI{4.3}{} to \SI{4.7}{\volt}: the trip current
        falls from \SI{143}{} to \SI{129}{\ampere} hot, and the stacked
        margin against normal operation from 1.33 to 1.16. This does not
        shorten the chain for a short present at turn-on, whose detection
        is set by the blanking alone; it helps only a fault that develops
        with the device already on.
  \item A larger zener is excluded. At \SI{5.1}{\volt} the lowest trip
        point, \SI{2.35}{\volt}, falls below the \SI{2.39}{\volt} of a
        hot maximum-resistance device carrying the \SI{60}{\ampere}
        overload: the protection would trip in normal operation.
\end{itemize}
%
With both changes the protection covers $t_\mathrm{SC} \geq
\SI{0.71}{\micro\second}$. Below that, desaturation detection is the wrong
principle for this device, and the answer is a different one: sensing
$\mathrm{d}i/\mathrm{d}t$ across the inductance between the power source
and the Kelvin source of the TO-247-4L, which sees a short within tens of
nanoseconds, or a lower gate voltage, which buys withstand time with
conduction loss and reopens C1.

\section{Current and voltage sensing}

\subsection{Shunt versus Hall}

\begin{table}[htbp]
\centering
\caption{Phase current sensing alternatives at \SI{30}{\Arms}.}
\label{tab:sensing}
\small
\begin{tabularx}{\textwidth}{lYY}
\toprule
 & \textbf{Shunt + isolated amplifier} & Hall-effect sensor \\
\midrule
Isolation & from the amplifier & intrinsic \\
Dissipation & \SI{0.45}{\watt} per phase & none \\
Offset drift & \textbf{very low} & significant \\
Linearity & \textbf{excellent} & degrades near full scale \\
External field sensitivity & \textbf{none} & yes \\
Bandwidth & \SI{280}{\kilo\hertz} & \SI{100}{}--\SI{200}{\kilo\hertz} \\
Cost & \textbf{low} & 3--5$\times$ \\
\bottomrule
\end{tabularx}
\end{table}

A \SI{0.5}{\milli\ohm} shunt with an \isoamp{} is selected. At
\SI{30}{\Arms} the shunt dissipates \SI{0.45}{\watt}, which is manageable,
and in exchange one obtains the stability and linearity of a metal
resistor. The choice would reverse above roughly \SI{100}{\ampere}, where
the shunt dissipation becomes the dominant problem.

\subsection{Where the shunt goes}

Three positions are possible: in series with the phase output, in the source
of the low-side device, or in the DC return.

The phase-series position is chosen, for three reasons in order of weight.
It is the only one that measures the phase current \emph{continuously} and in
both directions, which is what the current control loop needs; the low-side
position measures only during the low-side conduction interval, and that
window closes at the extremes of duty cycle. It stays out of the commutation
loop, whose inductance budget is already \SI{73}{\percent} consumed by the
device packages alone (\cref{sec:kelvin}). And the common-mode penalty it
incurs is covered by the chosen amplifier, though with less room than at
first sight: the shunt sits at the switching node, so the amplifier's input
side swings the full bus at up to \SI{70}{\volt\per\nano\second} over the
steepest nanosecond (\SI{31}{\volt\per\nano\second} between \SI{10}{} and
\SI{90}{\percent}) against a rated CMTI of
\SI{100}{\kilo\volt\per\micro\second} minimum~\cite{ti_amc1302}: a factor of
1.4 on the steepest value, 3.2 on the edge. This is the thinnest CMTI margin
on the board, and it belongs on the bench list.

What the phase position gives up is visibility of a shoot-through, whose
current never reaches the phase terminal. That fault is covered by
desaturation detection in \SI{200}{\nano\second}, which the shunt chain
could not have matched anyway.

\subsection{DC-link voltage}
\label{sec:dclinkv}

A resistive divider feeding an \isoampv{}~\cite{ti_amc1311}: six
\SI{680}{\kilo\ohm} resistors in the high side and \SI{8.2}{\kilo\ohm} in the
low side, a ratio of \num{498.6}:1 that puts the \SI{2}{\volt} full scale at
\SI{997}{\volt}. The high side is split because each 1206 part is rated
\SI{200}{\volt} and a single short-circuit failure must not place the bus on
the amplifier input.

\begin{finding}{The divider was at its rating before the voltage went up}
The first version used four \SI{1}{\mega\ohm} parts: \SI{200}{\volt} each
at \SI{800}{\volt}, which is \SI{100}{\percent} of the rating and not a
margin, and \SI{225}{\volt} at the \SI{900}{\volt} maximum continuous
voltage. Six parts carry \SI{150}{\volt} each at \SI{900}{\volt}
(\SI{75}{\percent}) and \SI{167}{\volt} at \SI{1000}{\volt}. The error was
invisible because ``four times \SI{200}{} is \SI{800}{}'' reads like a
sizing when it is a limit.
\end{finding}

\subsection{Overvoltage comparator}

The same measurement drives an on-board overvoltage trip at
$V_\mathrm{OV,trip} = \SI{950}{\volt}$ (\cref{tab:vhierarchy}), for the same
reason as the overload comparator of \cref{sec:analogif}: with the
controller off the board, a protection that exists only in firmware does not
exist on the board. At \SI{950}{\volt} the \isoampv{} differential output
is \SI{1.906}{\volt}, so its positive output sits \SI{0.953}{\volt} above
the common mode; one further comparator on the same reference pulls the
same fault line. The \SI{950}{\volt} is the level at which the switching
overshoot was verified (\cref{sec:kelvin}), which is what makes it the
right place to stop switching.

\section{Analogue interface to the controller}
\label{sec:analogif}

With the microcontroller outside the board boundary (\cref{sec:sysreq}), the
sensing chain does not end at an on-board ADC: it ends at a connector, and
the last stage is therefore an output driving a cable.

\subsection{No buffer}

The \isoamp{} has an output impedance below \SI{0.2}{\ohm} and a
short-circuit current of \SI{\pm 14}{\milli\ampere}~\cite{ti_amc1302}. One
or two metres of twisted cable present a hundred picofarads or so, which
that output drives without difficulty, and the receiving ADC input is high
impedance.

An operational-amplifier buffer would add its own offset and its own drift
to a chain whose entire justification against the Hall-effect alternative
(\cref{tab:sensing}) was low offset drift. It would buy drive capability
that is not needed and spend the reason for the choice. What remains is a
series resistor per line, which limits the current if an output is shorted
at the connector, isolates the cable capacitance from the output stage, and
forms the filter described next.

\subsection{Anti-aliasing: why the filter is at \texorpdfstring{\SI{106}{\kilo\hertz}}{106 kHz} and not at \texorpdfstring{\SI{10}{\kilo\hertz}}{10 kHz}}
\label{sec:aafilter}

The sampling rate is an assumption, and it is recorded as such in
\cref{app:assumptions}: $f_s = \fsw = \SI{20}{\kilo\hertz}$, once per
switching period, synchronous with the carrier. Nyquist is therefore
\SI{10}{\kilo\hertz}.

Placing the filter corner at \SI{10}{\kilo\hertz} would be a mistake. It
would add \SI{16}{\micro\second} of group delay to a current loop whose
own bandwidth is a few kilohertz, and it would throw away the
\SI{280}{\kilo\hertz} bandwidth that the isolated amplifier was chosen for.
The two chains that share this output have opposite requirements: the
control chain wants attenuation above Nyquist, and the protection
comparator wants bandwidth and no delay at all.

The conflict is resolved by \emph{tapping twice}. The overload comparator
takes the amplifier output directly, at full bandwidth; the filter sits
between that tap and the connector, so the protection path carries none of
its delay.

\begin{equation}
  f_c = \frac{1}{2\pi \cdot 2R \cdot C}
      = \frac{1}{2\pi \cdot \SI{99.8}{\ohm} \cdot \SI{15}{\nano\farad}}
      = \SI{106.3}{\kilo\hertz} ,
  \qquad
  t_d = \frac{1}{2\pi f_c} = \SI{1.50}{\micro\second} .
\end{equation}

with $R = \SI{49.9}{\ohm}$ in each of the two lines and $C$ connected
differentially between them. The corner erodes a \SI{2}{\kilo\hertz}
control bandwidth by \SI{-0.018}{\percent}, which is nothing, and a short from
either output to ground demands \SI{50}{\milli\ampere} from a stage that
self-limits at \SI{14}{\milli\ampere}~\cite{ti_amc1302}.

The switching ripple is removed by the sampling instant, as the next
subsection shows. The filter stops the wideband noise and
whatever the cable picks up, and \SI{106}{\kilo\hertz} does that job without
costing bandwidth or delay where either matters.

\subsection{Synchronous sampling, synchronised by the controller}
\label{sec:syncsampling}

The phase current carries \SI{2.463}{\ampere} of peak-to-peak switching
ripple, \SI{5.80}{\percent} of the fundamental peak
(\cref{fig:correnti}). Where the ripple is read decides whether that
appears in the measurement:

\begin{table}[htbp]
\centering
\caption{Reading error against the sampling instant.}
\label{tab:sampling}
\small
\begin{tabularx}{\textwidth}{lY}
\toprule
Sampling instant & Error in the current reading \\
\midrule
\textbf{Synchronous} with the carrier peak or valley
  & \textbf{zero by construction}: the ripple triangle crosses its mean there \\
\textbf{Asynchronous}, arbitrary instant
  & up to $\pm\SI{1.23}{\ampere}$, i.e.\ $\pm\SI{4.10}{\percent}$ of the
    \SI{30}{\Arms} rating; RMS value of the error \SI{0.711}{\ampere}
    (\SI{2.37}{\percent}) \\
\bottomrule
\end{tabularx}
\end{table}

Four percent on the feedback of a current loop is not acceptable, so
synchronous sampling is a requirement. It does \emph{not} follow that the
board should export a synchronisation line. The carrier does not exist on
this board: the board receives six already-modulated PWM signals. The
carrier lives in the controller, which therefore already owns the
synchronism; this is why every motor-control microcontroller
has its ADC trigger tied to the PWM timer. Sending a synchronisation signal
back would return information to where it came from, with the delay and the
jitter of the round trip added to it.

What the board \emph{does} have to publish is its own delay.

\begin{table}[htbp]
\centering
\caption{Delay from the phase current to the connector.}
\label{tab:chaindelay}
\begin{tabular}{lccl}
\toprule
Contribution & Typical & Maximum & Source \\
\midrule
\isoamp{}, input to output at \SI{50}{\percent}
  & \SI{1.6}{\micro\second} & \SI{2.1}{\micro\second} & \cite{ti_amc1302} \\
Anti-aliasing filter, $1/2\pi f_c$
  & \SI{1.50}{\micro\second} & \SI{1.50}{\micro\second} & \cref{sec:aafilter} \\
\midrule
\textbf{Total} & \textbf{\SI{3.1}{\micro\second}} & \textbf{\SI{3.6}{\micro\second}} & \\
\bottomrule
\end{tabular}
\end{table}

At the sampling instant the ripple crosses its mean, so its slope is at a
maximum: a delay is read along the steepest part of the waveform, which is
the worst case, and the one quoted here. The ripple slope is
\SI{0.0985}{\ampere\per\micro\second} and the fundamental adds
\SI{0.0533}{\ampere\per\micro\second} at its zero crossing, so
\SI{3.1}{\micro\second} displaces the reading by \SI{471}{\milli\ampere},
\SI{1.57}{\percent} of rating, and the worst case by
\SI{1.82}{\percent}.

This is a constant, known delay and not a random error: the controller
compensates it by advancing the ADC trigger, which is requirement R7 of
\cref{tab:sysreq}. The residual is trigger jitter, read on the same slope,
ripple and fundamental together (\SI{0.152}{\ampere\per\micro\second}):
staying below \SI{0.1}{\percent} of rating requires less than
\SI{198}{\nano\second} of it, which a hardware trigger tied to the PWM timer
meets with very large margin.

\begin{finding}{The filter costs almost as much delay as the amplifier}
The \SI{1.50}{\micro\second} of the anti-aliasing filter is nearly the
\SI{1.6}{\micro\second} of the isolated amplifier itself: half of the total
chain delay is added by two passive components. It remains the right choice
because that delay is constant and compensable whereas the noise it removes
is not, and because the protection path, which could not tolerate the
delay, taps upstream of it. The finding is recorded because the natural
assumption, that a passive filter two decades above the loop bandwidth is
free, is wrong here.
\end{finding}

\subsection{Overload comparator on the board}

\cref{tab:protection} assigns overload protection to a comparator on the
shunt chain. With the microcontroller outside the boundary that comparator
cannot be left to firmware, or the board would have no autonomous overload
protection at all: desaturation detection alone trips between
\SI{143}{} and \SI{259}{\ampere}, not at \SI{60}{\ampere}.

The comparator is therefore on the board, taking its input at the full-band
tap of \cref{sec:aafilter} and pulling the fault line low. Six comparators
are needed rather than two, because the amplifier output is differential
\emph{and} bipolar: at $\pm\SI{60}{\ampere}$ the differential output is
$\pm\SI{1.23}{\volt}$ about a \SI{1.44}{\volt} common
mode~\cite{ti_amc1302}, so the positive current raises OUTP to
\SI{2.06}{\volt} and the negative current raises OUTN to the same value.
Two thresholds per phase against one common reference, wire-ORed through
open-drain outputs, need three dual packages and a single voltage
reference.

The threshold is fixed while the amplifier common-mode output has a
tolerance, so common-mode drift adds directly to the trip point. For a
coarse trip at \SI{60}{\ampere} on a \SI{100}{\ampere} full scale that is
acceptable; this is a protection, not a measurement, and keeping it
separate from the control chain is deliberate.

\subsection{Connector}

The control connector is a \num{2}$\times$\num{20} IDC header. Row A carries
exactly twenty signals and row B is entirely ground, so every conductor in
the ribbon has a return beside it rather than a return shared with its
neighbour.

\begin{table}[htbp]
\centering
\caption{Row A of the control connector; row B, pins 2 to 40, is ground.}
\label{tab:connector}
\small
\begin{tabular}{llll}
\toprule
Pins & Signals & Count & Type \\
\midrule
A1--A11 & PWM H1, L1, H2, L2, H3, L3 & 6 & inputs, switching \\
A13, A15 & enable, discharge enable & 2 & inputs, static \\
A17, A19 & fault, ready & 2 & open-drain outputs, static \\
A21--A31 & phase currents 1--3, P and N & 6 & analogue, differential \\
A33, A35 & DC-link voltage, P and N & 2 & analogue, differential \\
A37, A39 & thermistors 1 and 2 & 2 & analogue \\
\midrule
 & & \textbf{20} & \\
\bottomrule
\end{tabular}
\end{table}

The four static lines separate the switching block from the analogue one,
so no analogue pin is adjacent to a switching edge. The \SI{12}{\volt}
supply (\SI{155}{\milli\ampere} by the budget of
\cref{tab:auxbudget}) leaves on a separate two-way connector, which is good
practice and also what makes the twenty signals fit exactly.

A ground conductor between the two halves
of a differential pair reduces crosstalk but decouples the pair, so the
differential coupling is worse than in a twisted pair. At $\pm\SI{2}{\volt}$,
filtered to \SI{106}{\kilo\hertz}, over a metre of cable, that is the right
trade; beyond two metres or in a harsher environment the answer would be a
connector with individually shielded pairs.

Each analogue line carries a low-capacitance bidirectional TVS to ground
downstream of the series resistor. Low capacitance matters: a
\SI{100}{\pico\farad} device across one line of a pair whose filter
capacitor is \SI{15}{\nano\farad} would unbalance the pair.

\section{DC-link precharge and discharge}
\label{sec:dclinkmgmt}

The DC-link is \SI{140}{} plus \SI{60}{\micro\farad} of film in the
\SI{1100}{\volt} class~\cite{vishay_mkp1848}, chosen in \cref{sec:dclink}.
The stored energy is $E = \tfrac12 C V^2 = \SI{64}{\joule}$ at
\SI{800}{\volt} and \SI{81}{\joule} at the \SI{900}{\volt} maximum
continuous voltage, enough to be dangerous and enough to weld a contactor if
switched in directly.

\subsection{Where the board stops}
\label{sec:scope}

The contactors and the precharge resistor are outside the board.
The converter receives an already-charged bus, and its high-voltage
connector carries five terminals: \texttt{DC+}, \texttt{DC$-$} and the three
phases. The precharge sizing is still given below, because whoever
integrates the system needs it.

The board's own contribution to precharge is the \emph{measurement}. The
main contactor should close when the DC-link has reached the
battery voltage, not when a timer expires: closing on a timer assumes the
precharge went as planned, and a partially charged bus welds the contactor
it was supposed to protect. The voltage that makes the decision possible is
the one produced by the isolated divider of \cref{sec:dclinkv}, reported to
the vehicle controller.

Sized as a system, the precharge uses a \SI{1}{\kilo\ohm} resistor, which
limits the inrush to \SI{0.9}{\ampere} at \SI{900}{\volt} and charges the
bus in about \SI{1}{\second} (five time constants of
\SI{0.2}{\second}), absorbing \SI{81}{\joule} as a pulse.

The interlock loop follows the same boundary. Its purpose is to command the
contactors open before a connector separates; with the contactors outside
the board, a local detector would raise an alarm the board cannot act on.
The only response available to the board, stopping the modulation, is one
the vehicle controller already commands through the enable line, and the
controller knows more about the rest of the vehicle than this board does. The board therefore carries the
loop through its high-voltage connector and does not monitor it.

\subsection{Active discharge}

Discharge must bring the bus below \SI{60}{\volt} within \SI{5}{\second} of
a crash signal~\cite{unr100,iso6469}. A purely passive bleeder sized for
that requirement from \SI{900}{\volt} would need \SI{9.2}{\kilo\ohm},
dissipating \SI{88}{\watt} continuously, more than three quarters of the
converter's own losses, spent doing nothing. An \emph{active} discharge is
therefore used: a \SI{1}{\kilo\ohm} resistor switched by a high-voltage
MOSFET, discharging from \SI{900}{\volt} to \SI{60}{\volt} in
\SI{0.54}{\second}. The resistor must absorb the whole stored energy as one
pulse, \SI{81}{\joule} at \SI{900}{\volt} and \SI{100}{\joule} at the
\SI{1000}{\volt} abnormal level: it is specified for
\SI{100}{\joule}, where the first version, sized for \SI{100}{\micro\farad}
at \SI{800}{\volt}, asked for \SI{32}{}.

The switch itself is undemanding thermally: in conduction it sits in
saturation, so the energy goes into the resistor and the device sees
milliwatts. Its constraint is blocking, because with no current
flowing the series resistor drops nothing and the full bus stands across it:
\SI{1200}{\volt} rating, \SI{1}{\ampere}, any on-resistance far below
\SI{1}{\kilo\ohm}. The E3M0040120K is used~\cite{wolfspeed_e3m0040}: the
automotive counterpart of the part first chosen, in the same package and
with the same drive as the power devices.

It is driven by a \emph{seventh instance of the gate-drive channel},
identical to the six that drive the power devices apart from desaturation
detection, which is disabled. A simpler optocoupler powered from the bus
itself was rejected for a specific reason: it stops working at low bus
voltage, which is exactly the end of the discharge it is supposed to
complete. A channel fed from the auxiliary rail holds the gate down to
\SI{0}{\volt} of bus.

\begin{finding}{The active discharge does not survive loss of the \SI{12}{\volt} rail}
The seventh channel is powered from the auxiliary rail. If a crash
interrupts the vehicle's \SI{12}{\volt} supply, the active
discharge does not run, and only the passive bleeder remains:
%
\begin{equation}
  t_{60} = R_\mathrm{bleed} C \ln\frac{900}{60}
  = \SI{224}{\kilo\ohm}\cdot\SI{200}{\micro\farad}\cdot 2.71
  = \SI{121}{\second},
\end{equation}
%
against a \SI{5}{\second} requirement: twenty-four times too slow.

This is the standard architecture, in which the bleeder is a backstop
against failure of the active circuit rather than a substitute for it. It is
recorded here as an explicit limit because \SI{121}{\second} against
\SI{5}{\second} is an unmet requirement in that condition.

A production design would remove the dependency in one of two ways: by
taking the discharge control power from the \emph{HV bus itself}, which is
the correct answer and costs an isolated supply rated for
\SI{1000}{\volt}; or by making the switch fail-safe closed, so that losing
control power \emph{causes} the discharge rather than preventing it, at the
price of turning every control fault into a bus discharge. Neither is
justified for a bench prototype whose \SI{12}{\volt} supply is a bench
lead; both would be mandatory in a vehicle.
\end{finding}

This is also why the auxiliary supply is taken from the vehicle's
\SI{12}{\volt} rail rather than derived from the HV bus: within the scope of
this board, the discharge circuit must stay powered while the bus is
collapsing.

\begin{finding}{The isolated supply cannot be fed from the raw vehicle rail}
The selected \dcdc{} accepts \SI{10.8}{}--\SI{13.2}{\volt}~\cite{murata_mgj2}.
A vehicle rail ranges from about \SI{9}{\volt} during cranking to
\SI{16}{\volt} while charging~\cite{iso16750}. Connecting the module directly
would destroy it at the first charging cycle. A pre-regulator is a mandatory
part of the design, and it was missing from the first bill of materials.
\end{finding}

\subsection{Auxiliary power budget}

The pre-regulator is sized by the budget of \cref{tab:auxbudget}, quoted in
full.

\begin{table}[htbp]
\centering
\caption{Auxiliary supply budget.}
\label{tab:auxbudget}
\begin{tabular}{lrrr}
\toprule
Load & Qty & Each & Total \\
\midrule
Gate driver, secondary side, \device{} & 6 & \SI{169.2}{\milli\watt} & \SI{1015.2}{\milli\watt} \\
Gate driver, secondary side, discharge  & 1 & \SI{117.6}{\milli\watt} & \SI{117.6}{\milli\watt} \\
\driver{}, primary side              & 7 & \SI{12.2}{\milli\watt}  & \SI{85.5}{\milli\watt} \\
\isoamp{}, high side                 & 3 & \SI{36.0}{\milli\watt}  & \SI{108.0}{\milli\watt} \\
\isoamp{}, low side                  & 3 & \SI{19.5}{\milli\watt}  & \SI{58.4}{\milli\watt} \\
\isoampv{}                           & 1 & \SI{40.0}{\milli\watt}  & \SI{40.0}{\milli\watt} \\
Interlock and fault logic            & 1 & \SI{50.0}{\milli\watt}  & \SI{50.0}{\milli\watt} \\
Overload comparators and reference   & 1 & \SI{10.0}{\milli\watt}  & \SI{10.0}{\milli\watt} \\
\midrule
\textbf{Loads}                       &   &   & \textbf{\SI{1484.7}{\milli\watt}} \\
Upstream of the isolated converters, $\eta = \num{0.80}$~\cite{murata_mgj2}
                                     &   &   & \SI{1855.9}{\milli\watt} \\
\textbf{Current from the \SI{12}{\volt} rail} & & & \textbf{\SI{155}{\milli\ampere}} \\
\bottomrule
\end{tabular}
\end{table}

Seven channels, not six: the seventh drives the active-discharge switch.
Each channel draws $(Q_g\fsw + I_\mathrm{Q})$ from the \SI{20}{\volt} gate
swing of \cref{sec:negbias}, with $I_\mathrm{Q} = \SI{4}{\milli\ampere}$ of
driver quiescent current~\cite{ti_ucc21750}: \SI{169.2}{\milli\watt} with
the \SI{223}{\nano\coulomb} of the \device{}, \SI{117.6}{\milli\watt} with
the \SI{94}{\nano\coulomb} of the discharge switch. The \SI{10}{\milli\watt}
allowed for the comparators is a deliberately generous round figure (seven
low-power comparators and a shunt reference biased at half a milliampere
come to about \SI{5}{\milli\watt}), and doubling it would still leave the
total under one percent.

Under two watts means the pre-regulator can be small. It is nevertheless
sized with a factor of four to cover the inrush at power-up, when seven
isolated modules and four sensing supplies charge their output capacitors
together, and for the same reason a \SI{220}{\micro\farad} bulk capacitor
sits immediately downstream of it. The part chosen at capture is an
\SI{8}{\watt} regulated module with a \SI{9}{}--\SI{18}{\volt}
input~\cite{traco_ten8} (D11, \cref{tab:captdecisions}); the four amplifier
supplies added at capture (D10) raise the current from the \SI{12}{\volt}
rail to about \SI{185}{\milli\ampere}, still a quarter of what it delivers.

\section{Hardware interlock}

Shoot-through must be impossible regardless of the software state. The
\driver{} has two command inputs, non-inverting and inverting, and the output
is active only when IN+ is high \emph{and} IN$-$ is low. Cross-connecting
them provides the interlock at no cost:

\begin{table}[htbp]
\centering
\caption{Cross-connected command inputs, leg $k$.}
\label{tab:interlock}
\begin{tabular}{lll}
\toprule
Driver & IN+ & IN$-$ \\
\midrule
High side & $\mathrm{PWM\_H}k$ & $\mathrm{PWM\_L}k$ \\
Low side  & $\mathrm{PWM\_L}k$ & $\mathrm{PWM\_H}k$ \\
\bottomrule
\end{tabular}
\end{table}

If both PWM signals are high, each driver sees IN+ = 1 and IN$-$ = 1 and
stays off. Both devices turn off: the fault state is safe by construction,
with no added component and no added delay. The dead time remains the task
of the microcontroller's PWM peripheral, which generates it precisely; the
interlock is an independent safety net, not the dead-time source.

The alternative of discrete logic gates would achieve the same with two more
packages and a few nanoseconds of delay. A single PWM line per leg with
hardware complementation would also work, but the dead time would then have
to come from an RC network whose threshold varies between
\num{0.3} and \num{0.7} of the supply, so that the same network gives a
delay anywhere within a factor of 3.4: too imprecise for a dead time whose
calculated minimum (\SI{166}{\nano\second}) and adopted value
(\SI{260}{\nano\second}) are only a factor of 1.57 apart.

\section{What the board requires from the system}
\label{sec:sysreq}

The microcontroller is outside the board boundary, for the same reason the
contactors, the precharge and the interlock supervisor are: \emph{the board
does not take in a function it cannot be the actuator of}. It receives six
modulated PWM signals and publishes measurements and a fault line; it does
not close a control loop.

The same boundary keeps discontinuous modulation at standstill out of the
sizing. It would lower the hottest junction at locked rotor by
\SI{17}{\celsius} (\cref{sec:lrdpwm}), but the modulation is generated by
the controller, outside this board, and hardware is not sized on a control
strategy the board cannot guarantee. A firmware update, a fallback mode after a fault, or a different
controller would remove the margin without anything on the board changing.
The margin is therefore written as a recommendation to the system, R13,
kept apart from the requirements the board depends on.

That boundary holds only if what crosses it is written down. The strongest
argument against it concerns the dead time: the \SI{260}{\nano\second} of
\cref{sec:deadtime} is generated by the PWM peripheral of a device that has
not been selected, and the whole dead-time analysis rests on it. The answer
is the same as for the other three boundaries (the vehicle controller
closes the main contactor on measured voltage, the vehicle supervises the
interlock loop, the controller generates the dead time), and it holds as
long as each is written as a requirement rather than left as a silent
assumption.

\begin{table}[htbp]
\centering
\caption{Requirements on the controlling system. R1--R12 and R14 are
mandatory: the board's verification depends on them. R13 is recommended: it
adds margin the verification does not count on.}
\label{tab:sysreq}
\small
\begin{tabularx}{\textwidth}{lYY}
\toprule
 & Requirement & Why \\
\midrule
R1 & Six PWM signals, complementary in pairs, \SI{3.3}{\volt} CMOS, with a
     dead time of at least \SI{260}{\nano\second} generated by the
     peripheral
   & \cref{sec:deadtime}: the calculated minimum is \SI{166}{\nano\second};
     the cross-connected interlock is a safety net, not the dead-time source \\
R2 & Enable low in every indeterminate state: reset, boot, watchdog
   & it is the only input that brings the board to its safe state \\
R3 & A desaturation fault is reset \emph{only} by holding enable low for
     more than \SI{1}{\micro\second}
   & the driver's reset and enable share one pin~\cite{ti_ucc21750} \\
R4 & Modulation is enabled only after the ready line goes high
   & all seven undervoltage lockouts must have released \\
R5 & The main contactor closes on the measured DC-link voltage, not on a
     timer
   & closing onto a half-charged bus welds the contactor it was meant to
     protect \\
R6 & Differential ADC inputs, sampled synchronously with the carrier peak
     or valley
   & \cref{sec:syncsampling}: asynchronous sampling costs
     \SI{4.10}{\percent} of reading error \\
R7 & The ADC trigger is advanced by \SI{3.1}{\micro\second}
   & \cref{tab:chaindelay}; uncompensated it costs \SI{1.57}{\percent} \\
R8 & Derating from the heatsink thermistor
   & the thermistor does not protect against transients; desaturation
     detection does \\
R9 & Supervision of the interlock loop and command of the contactors
   & the board carries the loop through without monitoring it \\
R10 & Galvanic isolation in the charging path, on-board in AC charging or
      in the station in DC charging
   & the contactors are \emph{closed} while charging; without the
     charger's transformer the overvoltage category applies again
     (\cref{sec:ovc}) \\
R11 & Source impedance seen by the DC-link at \SI{20}{\kilo\hertz}:
      $L \geq \SI{1.5}{\micro\henry}$ \emph{or} $R \geq \SI{50}{\milli\ohm}$
    & otherwise the source--capacitor resonance falls on the switching
      harmonics and the ripple exceeds the specification; the input filter
      of \cref{tab:emcfilter} meets it from inside the inverter with the
      \SI{12}{\micro\henry} leakage of its chokes \\
R12 & Maximum continuous battery voltage, charging and regenerative braking
      included: \SI{900}{\volt}; up to \SI{1000}{\volt} only with the bridge
      disabled and for less than \SI{30}{\minute} a day
    & \SI{900}{\volt} is the working voltage for creepage, device derating
      (\SI{75}{\percent}) and the DC-link; the film capacitor tolerates
      \num{1.15} times its rating for \SI{30}{\minute} a
      day~\cite{vishay_mkp1848} \\
\midrule
\multicolumn{3}{l}{\emph{Recommended}} \\
R13 & Below a few hertz of output frequency, a discontinuous zero sequence
      that clamps to the positive rail the phase carrying the largest
      positive current (DPWM-MAX)
    & at locked rotor the hottest junction falls from \SI{111.1}{} to
      \SI{93.9}{\celsius} and the converter loss from \SI{141}{} to
      \SI{96}{\watt}; in the simultaneous worst case, from \SI{142}{} to
      \SI{126}{\celsius} (\cref{sec:lrdpwm}). C1 is verified without it \\
\midrule
\multicolumn{3}{l}{\emph{Mandatory, added with the cooling rail}} \\
R14 & Coolant at the rail inlet: 50/50 water--glycol, at most
      \SI{65}{\celsius}, at least \SI{10}{\litre\per\minute}; the rail
      takes \SI{4.7}{\kilo\pascal}
    & the rail and the pad are sized on it (\cref{sec:cooling}); the
      worst case of \cref{tab:coolalt} is computed at \SI{70}{\celsius} \\
\bottomrule
\end{tabularx}
\end{table}

These fourteen lines are also written on the root sheet of the schematic,
with the recommended one marked as such.
A requirement that exists only in a report is one that the next engineer to
open the project will not read.
